\section{Distributed Massive MIMO Uplink System Model}
\label{sec:dmimo_ul_sysmodel}
We now consider the uplink of the same distributed (cell-free) massive \gls{mimo}
network as in \Cref{sec:dmimo_dl_sysmodel}. The $K$ users now transmit simultaneously on the same
time-frequency resources, each AP observes a superposition of all of them on its
$M$ antennas, and the per-AP observations are combined into one soft estimate per
user. Under \gls{tdd} operation the two directions occupy the same band within
one coherence block, so the propagation channels are the same $\vect{h}_{kl}[q]$
as in \eqref{eq:collective-channel} and the large-scale fading coefficients
$\beta_{kl}$ of \eqref{eq:beta-estimate} are unchanged. The downlink model
applies the array response as $\vect{h}_{k}[q]^{H}$, so reciprocity makes
$\vect{h}_{k}[q]$ itself the channel that carries the signal of user $k$ to the
$LM$ distributed antennas. 

User $k$ transmits $x_{k}[q]=\sqrt{p_{k}[q]}\,s_{k}[q]$ on subcarrier $q$, where
$s_{k}[q]\in\mathbb{C}$ is its unit-power data symbol,
$\expt{\abs{s_{k}[q]}^{2}}=1$, mutually independent across users, and
$p_{k}[q]\geq 0$ is the power it spends on that subcarrier. Each user is its own
transmitter and is subject to its own budget over the OFDM block,
\begin{equation}
\label{eq:ul-power-constraint}
\sum_{q=1}^{Q} p_{k}[q] = p_{k} \leq p_{\max},
\qquad k=1,\dots,K,
\end{equation}
with $p_{\max}$ the maximum transmit power of a user terminal. Since the power control below acts on the large-scale fading, which is flat over the band, the power is spread uniformly over the
subcarriers, $p_{k}[q]=p_{k}/Q$.

The signal received by AP $l$ on subcarrier $q$ is
\begin{equation}
\label{eq:ul-ap-received}
\vect{y}_{l}[q] = \sum_{k=1}^{K} \sqrt{p_{k}[q]}\,\vect{h}_{kl}[q]\, s_{k}[q]
+ \vect{n}_{l}[q] \inC{M},
\end{equation}
where $\vect{n}_{l}[q]\sim\cn{\vect{0}}{\sigma^{2}\mat{I}_{M}}$ is the receiver
noise, independent across APs and subcarriers and independent of the channels and
data. The noise power $\sigma^{2}$ is the same quantity as in
\eqref{eq:downlink-received}, but it is now referred to the AP receive chain, so
it is the one set by the noise figure of the LNAs and analog front end of
\Cref{sec:ana}. Stacking the $L$ observations gives the collective received vector
\begin{equation}
\label{eq:ul-collective-received}
\vect{y}[q] =
\begin{bmatrix} \vect{y}_{1}[q] \\ \vdots \\ \vect{y}_{L}[q] \end{bmatrix}
= \mat{H}[q]\,\mat{P}[q]^{1/2}\,\vect{s}[q] + \vect{n}[q] \inC{LM},
\end{equation}
with $\mat{H}[q]=[\,\vect{h}_{1}[q]\ \cdots\ \vect{h}_{K}[q]\,]\inC{LM\times K}$
the collective channel matrix already used in \Cref{sssec:centralized},
$\vect{s}[q]=[\,s_{1}[q]\ \cdots\ s_{K}[q]\,]^{\intercal}\inC{K}$,
$\vect{n}[q]\sim\cn{\vect{0}}{\sigma^{2}\mat{I}_{LM}}$, and
$\mat{P}[q]^{1/2}\inC{K\times K}$ the diagonal matrix with
$[\mat{P}[q]^{1/2}]_{k,k}=\sqrt{p_{k}[q]}$ and zero off-diagonal entries.

The data symbol of user $k$ is estimated by applying a receive combining vector
$\vect{v}_{k}[q]\inC{LM}$ to \eqref{eq:ul-collective-received}, which separates
into the desired term, the multi-user interference, and the noise,
\begin{align}
\label{eq:ul-estimate}
\hat{s}_{k}[q] ={}& \vect{v}_{k}[q]^{H}\vect{y}[q] \nonumber\\
={}& \underbrace{\sqrt{p_{k}[q]}\,\vect{v}_{k}[q]^{H}\vect{h}_{k}[q]\, s_{k}[q]}
_{\text{desired signal}} \nonumber\\
&+ \underbrace{\sum_{\substack{i=1 \\ i\neq k}}^{K}
\sqrt{p_{i}[q]}\,\vect{v}_{k}[q]^{H}\vect{h}_{i}[q]\, s_{i}[q]}
_{\text{multi-user interference}}
+ \underbrace{\vphantom{\sum_{i=1}^{K}}\vect{v}_{k}[q]^{H}\vect{n}[q]}
_{\text{noise}} .
\end{align}
Writing $\vect{v}_{k}[q]=[\,\vect{v}_{k1}[q]^{\intercal}\ \cdots\
\vect{v}_{kL}[q]^{\intercal}\,]^{\intercal}$ in blocks of $M$ entries, with
$\vect{v}_{kl}[q]\inC{M}$ the part applied to AP $l$, the estimate is the sum of
per-AP contributions,
\begin{equation}
\label{eq:ul-distributive}
\hat{s}_{k}[q] = \sum_{l=1}^{L} \vect{v}_{kl}[q]^{H}\vect{y}_{l}[q].
\end{equation}
Linear combining is therefore distributive over the APs, and this is the property
the functional splits of \Cref{ssec:split} exploit: an AP that holds its own
combining coefficients can form its partial sum locally and forward one scalar
stream per user instead of $M$ sample streams.

Treating the multi-user interference in \eqref{eq:ul-estimate} as noise, and
using that the noise term has variance
$\sigma^{2}\norm{\vect{v}_{k}[q]}^{2}$, the effective uplink \gls{sinr} of user
$k$ on subcarrier $q$ is
\begin{equation}
\label{eq:ul-sinr}
\mathrm{SINR}_{k}^{\mathrm{ul}}[q] =
\frac{p_{k}[q]\,\abs{\vect{v}_{k}[q]^{H}\vect{h}_{k}[q]}^{2}}
{\displaystyle\sum_{\substack{i=1 \\ i\neq k}}^{K}
p_{i}[q]\,\abs{\vect{v}_{k}[q]^{H}\vect{h}_{i}[q]}^{2}
+ \sigma^{2}\norm{\vect{v}_{k}[q]}^{2}} .
\end{equation}
With the effective channel known at the decoder, an achievable rate for user $k$ on subcarrier $q$ is
$\log_{2}\!\big(1+\mathrm{SINR}_{k}^{\mathrm{ul}}[q]\big)$, and inserting these
per-user, per-subcarrier SINRs into the ergodic rate expression of
\Cref{sec:rate} \eqref{eq:rate} gives the delivered per-user uplink rate
\begin{align}
\label{eq:ul-rate-user}
R_{\mathrm{UL},k} ={}& \bar x_{\mathrm{UL}}\,\tau_{\mathrm{UL}}\,
(1-\tau_{\mathrm{UL,sig}})\,\frac{\tilde B}{Q} \nonumber\\
&\times \expt{\sum_{q=1}^{Q}\log_{2}\!\left(1
+\mathrm{SINR}_{k}^{\mathrm{ul}}[q]\right)},
\end{align}
with $R_{\mathrm{UL}}=\sum_{k=1}^{K}R_{\mathrm{UL},k}$ the delivered uplink sum
rate. It drives the channel-decoder consumption of \Cref{sec:dig} and, through
\eqref{eq:fh_ul}, the traffic-dependent part of the fronthaul.


\subsection{Combiner and Power Control Designs}
\label{ssec:ul_designs}

As in the downlink, the design splits into the choice of the receive directions
$\{\vect{v}_{k}[q]\}$, which depends on how much the APs cooperate, and the
choice of the transmit powers $\{p_{k}\}$, which is a per-user allocation driven
by the large-scale fading. The two are treated in turn, and the directions are
recomputed every coherence block whereas the powers can be held fixed over many
blocks.

Throughout, and consistently with \Cref{ssec:designs}, we assume \emph{perfect}
\gls{csi}: every combiner is built from the true channels,
$\hat{\vect{h}}_{k}[q]=\vect{h}_{k}[q]$, so there is no estimation error and no
pilot contamination. The pilot overhead is nevertheless charged to the rate
through the frame parameters of \Cref{sec:frame}. In the downlink-only
evaluation of \Cref{ssec:designs} the uplink phase carried nothing but pilots,
$\tau_{\mathrm{UL,sig}}=1$, which leaves \eqref{eq:ul-rate-user} at zero. An
uplink that carries data requires $\tau_{\mathrm{UL,sig}}<1$, and the coherence
block then has to be split explicitly: with $\tau_{c}$ samples per block of which
$\tau_{p}$ are pilots, the signalling fraction of the frame is
$\tau_{\mathrm{UL}}\,\tau_{\mathrm{UL,sig}}=\tau_{p}/\tau_{c}$ and the uplink
data prelog of \eqref{eq:ul-rate-user} is
$\tau_{\mathrm{UL}}(1-\tau_{\mathrm{UL,sig}})=\tau_{\mathrm{UL}}-\tau_{p}/\tau_{c}$,
which requires $\tau_{p}/\tau_{c}\leq\tau_{\mathrm{UL}}$. Mutually orthogonal
pilots additionally require $\tau_{p}\geq K$; this is what makes the
no-contamination assumption above self-consistent, and it is the assumption under
which the perfect-\gls{csi} model is the idealisation of a real pilot phase rather
than a different system.


\subsubsection{Centralized Operation}
\label{sssec:ul_centralized}

In centralized operation the APs forward their received samples to the \gls{cpu},
which holds $\vect{y}[q]$ in full and combines over all $LM$ antennas jointly
\cite{demir2021foundations}. This is the highest level of cooperation, the
receive counterpart of the joint precoding of \Cref{sssec:centralized}, and it
corresponds to splits S1 and S2 of \Cref{ssec:split}. Minimizing the mean squared
error $\expt{\abs{s_{k}[q]-\hat{s}_{k}[q]}^{2}}$ over $\vect{v}_{k}[q]$ for the
model \eqref{eq:ul-collective-received} gives the \gls{mmse} combiner
\begin{equation}
\label{eq:ul-mmse-combiner}
\vect{v}_{k}[q] = \left(\mat{H}[q]\,\mat{P}[q]\,\mat{H}[q]^{H}
+ \sigma^{2}\mat{I}_{LM}\right)^{-1}\vect{h}_{k}[q],
\end{equation}
which requires inverting an $LM\times LM$ matrix. Applying the push-through
identity $(\mat{A}\mat{B}+\sigma^{2}\mat{I})^{-1}\mat{A}
=\mat{A}(\mat{B}\mat{A}+\sigma^{2}\mat{I})^{-1}$ with $\mat{A}=\mat{H}[q]\mat{P}[q]$
and $\mat{B}=\mat{H}[q]^{H}$, and discarding the diagonal factor that the scale
invariance of \eqref{eq:ul-sinr} makes irrelevant, the same directions are
obtained from a $K\times K$ inverse,
\begin{equation}
\label{eq:ul-centralized-rzf}
\mat{V}[q] = \mat{H}[q]\left(\mat{H}[q]^{H}\mat{H}[q]
+ \sigma^{2}\mat{P}[q]^{-1}\right)^{-1} \inC{LM\times K},
\end{equation}
whose column $k$ is $\vect{v}_{k}[q]$ up to a scalar. This is the uplink form of
the regularized family \eqref{eq:centralized-rzf}, with the heuristic loading
$\lambda\mat{I}_{K}$ replaced by the physically determined
$\sigma^{2}\mat{P}[q]^{-1}$; for equal powers $p_{k}[q]=p$ the two coincide at
$\lambda=\sigma^{2}/p$, which is the sense in which the downlink \gls{rzf}
precoder is the dual of \gls{mmse} combining. Letting the loading vanish, which
is the high-SNR limit $\sigma^{2}/p\to 0$, gives the \gls{zf} combiner
\begin{equation}
\label{eq:ul-zf-combiner}
\mat{V}[q] = \mat{H}[q]\big(\mat{H}[q]^{H}\mat{H}[q]\big)^{-1},
\end{equation}
for which $\vect{v}_{k}[q]^{H}\vect{h}_{i}[q]=0$ for all $i\neq k$, so all
inter-user interference is removed and \eqref{eq:ul-sinr} reduces to
$p_{k}[q]/(\sigma^{2}\norm{\vect{v}_{k}[q]}^{2})$. As in the downlink this needs
$K\leq LM$ for the Gram matrix to be invertible, and it disregards the noise, so
it is the scheme of choice only when the loading is negligible. Zero-forcing is
the scheme adopted in the power model of \Cref{sec:pwr_model} and its distributed
extension.

Because the combiner is scale-invariant, the centralized design carries no power
allocation of its own; the transmit powers are set at the users by
\eqref{eq:ul-fractional-power} below. The \gls{cpu} does need the powers
$\{p_{k}\}$ to build \eqref{eq:ul-centralized-rzf}, but these are large-scale
quantities refreshed once per large-scale fading realisation, not once per
coherence block.


\subsubsection{Local Operation and Large-Scale Fading Decoding}
\label{sssec:ul_local}

In local operation each AP combines its own observation with its own \gls{csi}
and forwards one scalar per user, so no instantaneous \gls{csi} crosses the
fronthaul. This is split S3 of \Cref{ssec:split}. AP $l$ forms
$\hat{s}_{kl}[q]=\vect{v}_{kl}[q]^{H}\vect{y}_{l}[q]$ using either \gls{mr}
combining, which points along the local channel,
\begin{equation}
\label{eq:ul-mr-combiner}
\vect{v}_{kl}[q] = \vect{h}_{kl}[q],
\end{equation}
or the local regularized combiner built from AP $l$'s own channel matrix
$\mat{E}_{l}[q]=[\,\vect{h}_{1l}[q]\ \cdots\ \vect{h}_{Kl}[q]\,]\inC{M\times K}$,
\begin{align}
\label{eq:ul-local-rzf}
\mat{V}_{l}[q] &= \big(\mat{E}_{l}[q]\mat{P}[q]\mat{E}_{l}[q]^{H}
+\sigma^{2}\mat{I}_{M}\big)^{-1}\mat{E}_{l}[q] \inC{M\times K},\\
 l&=1,\dots,L,
\end{align}
whose column $k$ is $\vect{v}_{kl}[q]$. This is \eqref{eq:local-rzf} with the
heuristic loading again replaced by its physical value, and it inverts an
$M\times M$ matrix because the per-AP array is small. An AP can suppress only the
interference it observes itself, using the $M$ degrees of freedom of its own
array, so the coherent cross-AP interference rejection of
\eqref{eq:ul-centralized-rzf} is lost; the loading is what keeps the inverse well
conditioned in the usual regime $M<K$.

What remains is how the \gls{cpu} fuses the $L$ local estimates. It applies one
weight per AP and user,
\begin{equation}
\label{eq:ul-lsfd}
\hat{s}_{k}[q] = \sum_{l=1}^{L} a_{kl}[q]^{*}\,\hat{s}_{kl}[q],
\end{equation}
collected into $\vect{a}_{k}[q]=[\,a_{k1}[q]\ \cdots\ a_{kL}[q]\,]^{\intercal}
\inC{L}$. By \eqref{eq:ul-distributive} this is exactly the collective combiner
$\vect{v}_{k}[q]$ whose $l$-th block is $a_{kl}[q]\vect{v}_{kl}[q]$, so
\eqref{eq:ul-sinr} continues to apply unchanged and the fusion rule is part of the
combiner rather than a separate stage. Two choices are used. Equal weighting,
$a_{kl}[q]=1$, simply adds the local estimates and needs nothing at the \gls{cpu}
beyond the streams themselves. \Gls{lsfd} instead optimizes the weights from the
channel statistics: collecting the per-AP effective gains
$g_{kil}[q]\triangleq\vect{v}_{kl}[q]^{H}\vect{h}_{il}[q]$ into
$\vect{g}_{ki}[q]=[\,g_{ki1}[q]\ \cdots\ g_{kiL}[q]\,]^{\intercal}\inC{L}$, and
writing $\mat{D}_{k}[q]$ for the diagonal matrix with entries
$\norm{\vect{v}_{kl}[q]}^{2}$ (diagonal because the noise is independent across
APs), the weights that maximize the resulting use-and-then-forget \gls{sinr} are
\begin{align}
\label{eq:ul-lsfd-weights}
\vect{a}_{k}[q] ={}& \Big(\sum_{i=1}^{K} p_{i}[q]\,
\expt{\vect{g}_{ki}[q]\vect{g}_{ki}[q]^{H}}
+ \sigma^{2}\expt{\mat{D}_{k}[q]}\Big)^{-1} \nonumber\\
&\times \sqrt{p_{k}[q]}\ \expt{\vect{g}_{kk}[q]},
\end{align}
where the expectations are over the small-scale fading for fixed user positions
\cite{demir2021foundations}.
\todo{verify \eqref{eq:ul-lsfd-weights} against the LSFD expression in
\cite{demir2021foundations} before it is used in the evaluation.}
These weights are statistical, not instantaneous, so they are computed once per
large-scale fading realisation and cost the fronthaul nothing per coherence
block. The three fusion rules therefore span three cooperation levels at the same
per-AP combiner: local combining at a single AP, equal-weight summation at the
\gls{cpu}, and statistically weighted summation. Only the centralized combiner of
\eqref{eq:ul-centralized-rzf} uses instantaneous \gls{csi} at the \gls{cpu}, and
the SE expressions of the different levels must not be mixed.

Two things should be flagged about how \eqref{eq:ul-sinr} is used at these lower
levels. When the fusion weights are statistical, the \gls{cpu} does not know the
instantaneous effective channel $\vect{v}_{k}[q]^{H}\vect{h}_{k}[q]$, and the
rigorous achievable rate is then the use-and-then-forget bound, in which the mean
effective channel is treated as the useful gain and its fluctuation as extra
interference. Evaluating \eqref{eq:ul-rate-user} with the instantaneous
\eqref{eq:ul-sinr}, as is done in the downlink of \Cref{sec:dmimo_dl_sysmodel},
assumes instead that the effective channel is known wherever the decoding
happens, and is therefore optimistic for local operation by the amount of
residual channel hardening.
\todo{decide whether the local-operation rate is reported with the instantaneous
SINR, for consistency with the downlink, or with the use-and-then-forget bound,
which is the achievable one; the gap grows as $M$ shrinks.}


\subsubsection{Uplink Power Control}
\label{sssec:ul_power}

The powers $\{p_{k}\}$ are set at the users from the large-scale fading and, by
\eqref{eq:ul-power-constraint}, are constrained one user at a time. Full power,
$p_{k}=p_{\max}$, is the natural reference, but it is the worst case for
fairness: a user close to an AP then swamps a distant one sharing the same
resources. Fractional power control weights the users by their aggregate
large-scale gain $\beta_{k}=\sum_{l=1}^{L}\beta_{kl}$,
\begin{equation}
\label{eq:ul-fractional-power}
p_{k} = p_{\max}\,
\frac{\beta_{k}^{\,v}}{\max_{i}\beta_{i}^{\,v}},
\qquad k=1,\dots,K,
\end{equation}
which satisfies $p_{k}\leq p_{\max}$ for either sign of the exponent $v$, since
the normalizing maximum is attained by the strongest user when $v>0$ and by the
weakest when $v<0$. The exponent plays the same throughput-versus-fairness role
as in \eqref{eq:centralized-fractional}, with $v=0$ recovering full power and
$v=-1$ inverting the channel statistically, so that every user arrives at the
network with the same average total gain $p_{k}\beta_{k}$. The normalization by a
maximum rather than by a sum is what differs from the downlink rule
\eqref{eq:centralized-fractional}: there a shared budget had to be divided among
the users, whereas here each user has its own budget and the only role of the
denominator is to keep the strongest user within it. As a result the uplink has
no counterpart of the common rescaling \eqref{eq:power-normalization}, and no
feasibility question arises.

The uplink transmit power is spent at the users and therefore does not appear in
the network consumption $P_{\mathrm{net}}$ of \eqref{eq:pnet}, which counts the
APs, the fronthaul, and the \gls{cpu} only. The asymmetry this creates in the
power model is worth stating plainly: in the downlink phase the APs pay for the
PAs of \Cref{ssec:dmimo_pa}, whose consumption scales with the radiated power,
whereas in the uplink phase they pay only for the receive front end
\eqref{eq:ana_ul_ap} and the digital processing \eqref{eq:dig_ap}, both of which
are essentially independent of $p_{\max}$. Uplink power control therefore acts on
the delivered rate \eqref{eq:ul-rate-user} and on nothing else in
\eqref{eq:ee_net}, which makes it a pure spectral-efficiency lever rather than an
energy one, and it is the reason the uplink and downlink phases of the same frame
have very different consumption profiles.


\subsection{Complexity of combiners}
\todo{compute complexity of the combiners so this can be used in the power model;
under TDD reciprocity the centralized combiner reuses the matrix already
inverted for \eqref{eq:gops_cen}, so only the application cost $2K'N$ should be
charged in the uplink, whereas \gls{lsfd} adds the per-user $L\times L$ solve of
\eqref{eq:ul-lsfd-weights} once per large-scale fading realisation.}
